\section{Результаты II: механизм}
\label{sec:mechanism}

Раздел~\ref{sec:capability} показывает, что отброшенная недавняя история
обходится в решаемость задач. Этот раздел показывает \emph{как} именно
платится эта цена — по вызовам инструментов тех же самых запусков.

\subsection{Агент сам выкупает то, что вычеркнул промпт}

Запуск с одним лишь состоянием проваливается не потому, что забывает план.
Он проваливается потому, что заново выводит то, что уже успел
пронаблюдать. Разбираем по назначению каждую команду оболочки из 472
полностью трассированных запусков и считаем команды, повторяющие более
ранний идентичный вызов\live:

\begin{center}
\begin{tabular}{@{}lrrr@{}}
\toprule
Условие & Команд на запуск & Доля осмотра & \textbf{Повторный осмотр} \\
\midrule
\Control{} (полная история) & 95  & 0.23 & \textbf{0.20} \\
\StateOnly{}, решено      & 163 & 0.52 & 0.45 \\
\StateOnly{}, провалено      & 278 & 0.67 & \textbf{0.58} \\
\bottomrule
\end{tabular}
\end{center}

У провалившегося запуска с одним состоянием \SI{58}{\percent} всех команд
уходит на повтор уже выполненного осмотра — против \SI{20}{\percent} у
контроля. На отдельной матрице тот же эффект виден и по другому показателю:
переходы по указателям к файлам случаются 5.19 раза за запуск против 0.42 у
контроля — примерно в двенадцать раз чаще\live.

Блок состояния ограничивает \emph{промпт}, но не ограничивает
\emph{объём информации, которого требует задача}, и разницу агент
покрывает через цикл вызовов инструментов — именно там, куда не смотрит ни
одна метрика размера промпта. Это и есть тезис статьи, только
сформулированный как измерение, а не как довод.

\subsection{Прочитано вручную, одна траектория}

Одна траектория, изученная напрямую, а не через агрегат\live: 72 вызова
инструмента, из них \texttt{input.tex} прочитан 22 раза (двумя разными
написаниями пути), второй файл — 19 раз, документ перекомпилирован пять раз
и каждый раз выдаёт те же семь предупреждений на тех же номерах строк.
Ходы 22--33 — это почти сплошь чередующиеся повторные чтения двух файлов,
чьи байты ни разу не менялись.

Агрегат может произвести и сломанный метод подсчёта; здесь всё было наоборот
— сначала прочитали вручную, а метрику написали потом, чтобы воспроизвести
уже увиденное механически.

\subsection{Что на самом деле лежало в состоянии}

Смотрим на финальное состояние 251 запуска с одним состоянием, отдельно по
исходу\live:

\begin{center}
\begin{tabular}{@{}lrr@{}}
\toprule
Маркер в финальном состоянии & Решено & Провалено \\
\midrule
Символов (медиана)      & 1218 & 963 \\
Путей к файлам               & 9    & 5 \\
Точных версий    & 1    & \textbf{0} \\
Текста ошибки или кода завершения  & 0    & \textbf{0} \\
Процитированных команд          & 1    & \textbf{0} \\
\bottomrule
\end{tabular}
\end{center}

В состоянии провалившегося запуска нет ни одной точной версии, ни одного
текста ошибки, ни одной команды — то есть именно тех фактов, которые агент
затем заново добывает, повторно запуская ту же команду. Схеме попросту
некуда их записать: её поля — это прозаические контейнеры (\texttt{goal},
\texttt{facts\_established}, \texttt{work\_done}, \texttt{work\_remaining},
\texttt{next\_step}, \texttt{current\_blocker}).

Связь идёт в ожидаемую сторону\live: запуски, чьё состояние фиксирует не
менее восьми путей к файлам, решаются в 0.50 случаев против 0.24 у тех, что
фиксируют меньше; запуски с записанной точной версией — в 0.47 против 0.32;
запуски с записанной ошибкой или кодом завершения — в 0.50 против 0.33. Это
корреляция (запуск, продвинувшийся дальше, просто имел что записывать), но
это единственное направление, куда указывают данные, и оно и мотивирует
исследование композиции в \S\ref{sec:composition}.

\subsection{Почему исходная оценка не могла этого увидеть}

Повторное выяснение тратится дополнительными шагами по инициативе самого
агента. Генератор оригинала фиксирует число шагов на эпизод
(\S\ref{sec:repro}), так что механизму, измеренному в этом разделе, просто
негде проявиться. Именно в этом смысле синтетический результат не ошибочен,
а неприменим: бенчмарк убирает саму переменную, через которую идёт плата.
